\section{Comparação com a literatura}
\label{sec_comparacao_literatura}

A modelagem da dinâmica orbital relativa adotada nesta dissertação fundamenta-se nas equações de \gls{hcw}, amplamente utilizadas na literatura para descrever o movimento relativo entre espaçonaves em órbitas próximas.
Conforme discutido na Seção~\ref{sec_modelagem_orbital}, essa formulação linear permite representar de maneira adequada o comportamento relativo durante operações de proximidade, particularmente quando a distância entre os veículos é pequena em comparação ao raio orbital.
Assim, de forma consistente com os trabalhos apresentados por \cite{clohessy1960terminal} e com a interpretação discutida por \cite{fehse2003}, a modelagem adotada neste estudo mantém as mesmas hipóteses de linearização da dinâmica gravitacional.

As estratégias de controle utilizadas diferem das abordagens mais elaboradas frequentemente discutidas na literatura. 
Enquanto diversos estudos exploram técnicas de controle não linear, adaptativo ou robusto para lidar com incertezas e perturbações em operações de proximidade \cite{slotine1991applied}, nesta dissertação optou-se pela utilização de leis de controle do tipo proporcional--derivativa.
Essa escolha foi aplicada de forma consistente aos diferentes subsistemas considerados no modelo: controle de atitude da espaçonave, controle das juntas do manipulador robótico e controle da dinâmica translacional relativa baseada nas equações de \gls{hcw}. Dessa forma, a abordagem adotada privilegia a simplicidade de implementação e a interpretação clara do comportamento dinâmico do sistema, mantendo coerência com a aplicação de controladores \gls{pd} em sistemas aproximadamente lineares descritos por essas equações, conforme discutido por \cite{qi2022variable}.
Assim, diferentemente de trabalhos que concentram sua contribuição em leis de controle avançadas, o foco desta dissertação está na análise da interação dinâmica entre os subsistemas e na avaliação do desempenho obtido com uma estratégia de controle aplicada ao problema de proximidade orbital com manipuladores robóticos.

No modelo desenvolvido, o manipulador robótico é considerado como parte integrante do sistema dinâmico formado pela espaçonave servidora e pelo alvo durante operações de proximidade.
Trabalhos como o de \cite{papadopoulos1994free} investigam de forma geral a dinâmica de sistemas robóticos espaciais com base livre e o acoplamento entre manipulador e espaçonave.
Nesta dissertação, essa base conceitual é utilizada para representar o manipulador no cenário específico de operações de proximidade e atracação.
Além disso, estudos voltados à atracação de alvos não cooperativos, como o de \cite{nishida2011strategy}, enfatizam os desafios associados à transferência de momento e às incertezas nos parâmetros inerciais durante o contato.
Em comparação, o manipulador é aqui modelado no contexto da interação dinâmica com o sistema orbital relativo, permitindo analisar sua atuação em conjunto com a dinâmica da espaçonave servidora.
Dessa forma, diferentemente de trabalhos que se concentram no planejamento detalhado da atracação ou em estratégias específicas de controle do manipulador, o enfoque está na integração do manipulador robótico ao modelo dinâmico do sistema utilizado para estudar operações de serviço em órbita, conforme discutido em revisões como a de \cite{flores2014review}.

Por fim, a dinâmica acoplada é considerada no contexto da interação entre a dinâmica orbital relativa, a atitude da espaçonave servidora e o movimento do manipulador robótico.
Estudos experimentais anteriores, como os conduzidos com o satélite ETS-VI \cite{adachi1999}, tiveram como objetivo validar modelos dinâmicos de atitude por meio de dados obtidos em órbita, enquanto missões como o ETS-VII \cite{yoshida2003} permitiram investigar diretamente os efeitos de reação gerados por manipuladores robóticos instalados em plataformas com base livre.
Em comparação, esta dissertação não se concentra na validação experimental em voo, mas na modelagem e análise do comportamento dinâmico resultante da interação entre os diferentes subsistemas envolvidos.
Além disso, enquanto alguns trabalhos abordam o acoplamento entre posição relativa e atitude no contexto de operações de aproximação e acoplamento, como no controle integrado proposto por Wei et al.~\cite{wei2011}, o enfoque deste estudo está na análise do acoplamento dinâmico do sistema completo — incluindo dinâmica orbital relativa, atitude da espaçonave e movimento do manipulador — dentro de um modelo unificado utilizado para avaliar operações de proximidade e atracação em serviço em órbita.